\documentclass[10pt,a4paper]{amsart}
\usepackage[margin=19mm]{geometry}
\usepackage{amsmath,amssymb,mathtools}
\usepackage[scr=rsfs,cal=euler]{mathalfa}
\usepackage{xfrac}
\usepackage[unicode,hidelinks,bookmarks=false]{hyperref}
\newcommand{\ie}{i.e.\ }
\newcommand{\cf}{cf.\ }
\newcommand{\resp}{respectively\ }
\newcommand{\Subsection}{\subsection}
\providecommand{\nobreakdash}{\nobreak\hbox{-}}
\setlength{\emergencystretch}{2em}
\multlinegap0pt
\usepackage{amssymb}
\usepackage{xcolor}
\newtheorem{proposition}{Proposition}
\newcommand{\FS}[0]{\mathsf{FS}}
\newcommand{\NN}[0]{\mathbb{N}}
\newcommand{\RCA}[0]{\mathsf{RCA}}
\newcommand{\W}{\mathrm{W}}
\title{Ramsey-like theorems for the Schreier barrier}
\author{Lorenzo Carlucci, Oriola Gjetaj, Quentin Le Hou\'erou, Ludovic Levy Patey}
\date{Source: arXiv:2412.11598v3 (CC BY 4.0)}
\begin{document}
\maketitle

\noindent\emph{Source and licence.} Ramsey-like theorems for the Schreier barrier. Version arXiv:2412.11598v3 (CC BY 4.0), distributed by arXiv under the Creative Commons Attribution 4.0 International licence, http://creativecommons.org/licenses/by/4.0/, as recorded in the arXiv licence metadata for this exact version and shown on the abstract page https://arxiv.org/abs/2412.11598v3. Full licence notice: \texttt{licenses/arxiv-2412.11598-CC-BY-4.0.txt}.\par\smallskip

\noindent\textit{Editorial scope.} Counted unit: the article's proposition proposition (source0.tex lines 585--587). The article's complete original proof of it is delivered with it (source0.tex lines 589--608, one \texttt{proof} environment of 20 lines). No mathematics is written, repaired or re-derived: the statement and the proof are the article's own lines, in the article's own notation, and no retained line is substituted or re-flowed.  No further source-local context is needed: every cross-reference of every retained line resolves inside the two delivered spans.  Nothing was omitted from any retained span, so the package carries no omission record.  No retained line contains a citation, so the wrapper carries no bibliography and no bibliography entry.  No retained line contains a figure environment, an image-inclusion command or drawing code, so no image is delivered and the package declares no asset dependency.  Editorial formatting only, in addition: an AMS \texttt{amsart} preamble that restates the article's own declarations and macros used by the retained lines, namely the environments \texttt{proposition} on separate counters, together with the standard packages those lines need.  The retained environments print in this document as Proposition 1 in order of appearance, whereas the article prints the same items with the numbering of its own full document.  The complete native source of the article as distributed in the arXiv e-print for this exact version is bundled unchanged as \texttt{sources/170182/source0.tex}.
\begin{proposition}
$\FS^{!\omega} \geq_{\W} \forall n \FS^n$ and $\RCA_0\vdash \FS^{!\omega} \to \forall n \FS^n$
\end{proposition}

\begin{proof}
Let $f:[\NN]^n\to \NN$. Define $g:[\NN]^{!\omega} \to \NN$ as follows. For $s = \{s_0, s_1, \dots, s_{s_0}\}$ we define 
$g(s) = f(s_1, \dots, s_n)$ if $s_0 \geq n$, else $0$. 

Now let $M= \{ m_0 < m_1 < \dots\}$ be an infinite set such that for all $s\in [M]^{!\omega}$, if $g(s) \in M$ then 
$g(s) \in s$, as given by $\FS^{!\omega}$.
Let $m_j$ be the least element of $M$ larger than or equal to $n$. Let 
$$M_n=M \cap [m_j + 1, \infty) \setminus \{ m_{2i} \;:\; i > j \}.$$ 
We claim that $M_n$ is free for $f$. Let $\{x_1, \dots, x_n\}\subseteq M_n$. 
If $f(x_1, \dots, x_n) \notin M_n$ we have nothing to prove. Suppose then that
$f(x_1, \dots, x_n) \in M_n$. Then $f(x_1, \dots, x_n) \in M$. 
Let $s$ be a canonically chosen exactly $\omega$-large set $\{m_j\}\cup \{ x_1, \dots, x_n\} \cup T$
extending $\{x_1, \dots, x_n\}$ with minimum 
element $m_j$, followed by $\{x_1, \dots, x_n\}$, followed by a tail $T$ consisting of the first $m_j - n$ elements of $M$ with even index
beyond $x_n$. 
By definition of $g$, $g(s) = f(x_1, \dots, x_n)$. Since $g(s) \in M$ and $M$ is free for $g$, 
it must be the case that $g(s) \in s=\{m_j\}\cup \{ x_1, \dots, x_n\} \cup T$. Then, since $g(s)$ is also in $M_n$ by assumption, 
and $(\{m_j\}\cup T) \cap M_n = \emptyset$ by construction, 
it must be the case that $g(s) \in \{x_1, \dots, x_n\}$. This proves the claim. 
\end{proof}

\end{document}
